\section*{Solutions to Homework 6}

\textbf{Part A.}

\begin{lstlisting}
$ bash codeHints/L06/setup_hw6.sh
$ cd practice
$ git clone hw6_server.git hw6
$ cd hw6
$ git remote -v                      # origin ... (fetch), origin ... (push)
$ git branch -a
* main
  remotes/origin/HEAD -> origin/main
  remotes/origin/main
$ echo "my work" > mine.txt
$ git add mine.txt
$ git commit -m "Add my work"
$ git status                         # "ahead of 'origin/main' by 1 commit"
$ printf 'remote=origin\nahead=1\nno-merge=fetch\n' > ../hw6_answers.txt
$ cd ../..                           # back to the course folder
$ bash codeHints/L06/friend.sh
$ cd practice/hw6
$ git push                           # rejected (fetch first)
$ git config pull.rebase false
$ git pull --no-edit                 # fetch + merge: one merge commit
$ git log --oneline --graph
$ git push
$ git switch -c feature
$ echo "feature" > feature.txt
$ git add feature.txt
$ git commit -m "Add feature file"
$ git push -u origin feature
\end{lstlisting}

\textbf{6.1--6.2.} \code{git clone} created \code{hw6} with the full history and a remote named \code{origin} that points at the server folder. The remote's name is \code{origin} (Git chooses it for a clone), hence \code{remote=origin}.

\textbf{6.3.} Your branch has one commit that \code{origin/main} lacks, so \code{git status} says ``ahead by 1 commit'' and \code{ahead=1}. Nothing is on the server yet.

\textbf{6.4.} The friend's commit reached the server after you cloned, so the two copies diverged and \code{git push} is rejected with ``fetch first''. Plain \code{git pull} on a recent Git stops with ``Need to specify how to reconcile divergent branches'' (after already fetching); \code{git config pull.rebase false} chooses merge. The merge commit is the \emph{one} merge commit the checker expects. (\code{git fetch} followed by \code{git merge origin/main} gives the same result.)

\textbf{6.5.} After the merge your \code{main} is ahead of \code{origin/main}, and a push is now a fast-forward for the server, so it is accepted.

\textbf{6.6.} \code{-u} records that local \code{feature} tracks \code{origin/feature} (the checker reads \code{branch.feature.remote} and \code{branch.feature.merge} from your repository).

\textbf{6.7.} \code{fetch}: it downloads and moves \code{origin/main} but changes neither your branches nor your files. \code{git pull} also merges.

\textbf{Part B.}

\begin{lstlisting}
$ cd practice
$ mkdir hw6_github
$ cd hw6_github
$ git init -b main
$ echo "# Hello GitHub" > README.md
$ git add README.md
$ git commit -m "Add README"
$ git remote add origin git@github.com:YOURNAME/hello-github.git
$ git push -u origin main
        (now edit README.md on github.com, "Commit changes")
$ git pull
$ git log --oneline                  # two commits
\end{lstlisting}

Replace \code{YOURNAME} by your GitHub user name. After the browser edit, GitHub is one commit ahead of you, so \code{git pull} brings it in as a fast-forward and your local \code{main} again equals GitHub's.

\Needspace{8\baselineskip}
\textbf{Common mistakes and fixes.}
\begin{itemize}
  \item Running \code{friend.sh} before cloning, or pushing before 6.4. Then there is no divergence and no merge commit, so the check ``exactly one merge commit'' fails. Run \code{setup\_hw6.sh --reset} and follow the order.
  \item \code{Permission denied (publickey)} when pushing to GitHub: the key is not added to your account, or the remote uses the wrong address style. Test with \code{ssh -T git@github.com}; check \code{git remote -v}.
  \item Creating the GitHub repository \emph{with} a README, then pushing: rejected, because GitHub has a commit you lack. Either create it empty (best) or run \code{git pull origin main --allow-unrelated-histories}, resolve any conflict, and push.
  \item Using \code{git push --force} to ``fix'' a rejection: not allowed here, it erases others' work.
\end{itemize}
